\documentclass[a4paper,12pt]{article}

\usepackage[T2A]{fontenc}
\usepackage[utf8]{inputenc}
\usepackage[russian]{babel}
\usepackage{amsmath,amssymb,mathtools}
\usepackage{helvet}
\renewcommand{\familydefault}{\sfdefault}
\usepackage{icomma}
\usepackage[a4paper,margin=19mm,headheight=15pt]{geometry}
\usepackage{microtype}
\usepackage{tikz}
\usetikzlibrary{calc,arrows.meta,positioning,shapes.geometric}
\usepackage{tkz-euclide}
\usepackage{tabularx,booktabs}
\usepackage{enumitem}
\usepackage{hyperref}
\usepackage{parskip}
\usepackage{fancyhdr}
\usepackage{setspace}
\usepackage{xcolor}

\definecolor{accent}{HTML}{008080}
\definecolor{darkaccent}{HTML}{004D4D}
\definecolor{textgray}{HTML}{333333}

\hypersetup{hidelinks}
\onehalfspacing
\setlength{\parindent}{0pt}
\setlength{\emergencystretch}{2em}
\renewcommand{\arraystretch}{1.28}

\pagestyle{fancy}
\fancyhf{}
\fancyhead[L]{\small\textcolor{darkaccent}{Векторы на координатной плоскости}}
\fancyhead[R]{\small Матвей Горбачев}
\fancyfoot[C]{\small\thepage}
\renewcommand{\headrulewidth}{0.4pt}
\renewcommand{\footrulewidth}{0pt}

\tkzSetUpLine[line width=1pt,color=accent]
\tkzSetUpPoint[color=accent,fill=accent,size=3]
\tkzSetUpLabel[color=black,font=\small]

\tikzset{
  flowstart/.style={
    ellipse,
    draw=darkaccent,
    line width=0.9pt,
    align=center,
    minimum width=42mm,
    minimum height=10mm,
    inner sep=2.5mm,
    font=\small
  },
  flowaction/.style={
    rounded corners=2mm,
    draw=accent,
    line width=0.9pt,
    align=center,
    text width=68mm,
    minimum height=11mm,
    inner sep=2.7mm,
    font=\small
  },
  flowdecision/.style={
    diamond,
    aspect=2.3,
    draw=darkaccent,
    line width=0.9pt,
    align=center,
    text width=42mm,
    inner sep=1.3mm,
    font=\small
  },
  flowarrow/.style={
    -{Latex[length=3mm,width=2mm]},
    line width=0.9pt,
    draw=textgray
  }
}

\newcommand{\vect}[1]{\vec{#1}}
\newcommand{\checkitem}{\item[$\square$]}
\newcommand{\sectionrule}{%
  \par\vspace{1mm}%
  {\color{accent!60}\hrule height 0.7pt}%
  \vspace{3mm}%
}

\begin{document}

% =========================================================
% Титульный лист
% =========================================================

\begin{titlepage}
\thispagestyle{empty}
\centering

\vspace*{0.15\textheight}

{\Large\textcolor{darkaccent}{Чек-лист}\par}

\vspace{7mm}

{\Huge\bfseries Векторы на координатной плоскости\par}

\vspace{4mm}

{\Large
Координаты, длина, операции\\
и скалярное произведение
\par}



\vfill

{\normalsize \today\par}

\vspace{8mm}

{\normalsize
Лёвин Артём Александрович\\
эксклюзивно для Матвея Горбачева
\par}

\begin{tikzpicture}[remember picture,overlay]
\draw[
  accent!45,
  line width=1.5pt
]
  ($(current page.south west)+(0,18mm)$)
  .. controls
  ($(current page.south west)+(55mm,37mm)$)
  and
  ($(current page.south east)+(-58mm,7mm)$)
  ..
  ($(current page.south east)+(0,23mm)$);
\end{tikzpicture}

\end{titlepage}

% =========================================================
% Основной чек-лист
% =========================================================

\section*{1. Координаты вектора}
\sectionrule

Пусть
\[
A(x_A;y_A)
\]
--- начало вектора, а
\[
B(x_B;y_B)
\]
--- его конец. Тогда

\[
\boxed{
\overrightarrow{AB}
=
(x_B-x_A;\,y_B-y_A)
}
\]

\begin{itemize}[leftmargin=7mm,itemsep=2mm]

\checkitem
Всегда используйте порядок
\[
\text{координаты конца}-\text{координаты начала}.
\]

\checkitem
Координата по \(x\) показывает горизонтальное смещение:
вправо --- положительное, влево --- отрицательное.

\checkitem
Координата по \(y\) показывает вертикальное смещение:
вверх --- положительное, вниз --- отрицательное.

\checkitem
Если начало вектора совпадает с началом координат,
координаты вектора совпадают с координатами его конца.

\end{itemize}

\textbf{Пример.}
Пусть
\[
A(3;2),\qquad B(9;3).
\]
Тогда
\[
\overrightarrow{AB}
=
(9-3;\,3-2)
=
(6;1).
\]
Это означает смещение на \(6\) единиц вправо и на \(1\) единицу вверх.

\section*{2. Длина вектора}
\sectionrule

Если
\[
\vect a=(x_a;y_a),
\]
то
\[
\boxed{|\vect a|=\sqrt{x_a^2+y_a^2}}
\]
и, следовательно,
\[
\boxed{|\vect a|^2=x_a^2+y_a^2}.
\]

Если в задаче требуется именно \textbf{квадрат длины}, извлекать квадратный корень не нужно.
Например, для
\[
\vect a=(6;2)
\]
получаем
\[
|\vect a|^2=6^2+2^2=40.
\]

\section*{3. Основные операции с векторами}
\sectionrule

Для
\[
\vect a=(x_a;y_a),
\qquad
\vect b=(x_b;y_b)
\]
используются следующие правила.

\begin{table}[h!]
\centering
\begin{tabularx}{\linewidth}{@{}>{\bfseries}p{0.31\linewidth}X@{}}
\toprule
Операция & Формула \\
\midrule
Умножение на число
&
\(k\vect a=(kx_a;ky_a)\)
\\[1.5mm]
Сложение
&
\(\vect a+\vect b=(x_a+x_b;\,y_a+y_b)\)
\\[1.5mm]
Вычитание
&
\(\vect a-\vect b=(x_a-x_b;\,y_a-y_b)\)
\\[1.5mm]
Противоположный вектор
&
\(-\vect b=(-x_b;-y_b)\)
\\
\bottomrule
\end{tabularx}
\end{table}

Вычитание удобно понимать через сложение:
\[
\boxed{\vect a-\vect b=\vect a+(-\vect b)}.
\]
Противоположный вектор имеет ту же длину, но направлен в противоположную сторону:
\[
\overrightarrow{BA}=-\overrightarrow{AB}.
\]

\section*{4. Геометрическое сложение векторов}
\sectionrule

При сложении действует \textbf{правило треугольника}:
\[
\boxed{
\overrightarrow{AB}+\overrightarrow{BC}=\overrightarrow{AC}
}.
\]

Чтобы сложить два произвольных вектора геометрически, второй вектор можно параллельно перенести так,
чтобы его начало совпало с концом первого. Сумма соединяет начало первого вектора с концом перенесённого второго.

Для вычитания сначала замените
\[
\vect a-\vect b
\]
на
\[
\vect a+(-\vect b),
\]
после чего примените обычное правило сложения.

\textbf{Важно различать}
\[
|\vect a|+|\vect b|
\qquad\text{и}\qquad
|\vect a+\vect b|.
\]
В общем случае эти величины не равны.

\section*{5. Скалярное произведение и угол между векторами}
\sectionrule

Скалярное произведение двух векторов --- \textbf{число}. Его можно вычислять двумя способами.

\subsection*{Через координаты}

\[
\boxed{
\vect a\cdot\vect b
=
x_ax_b+y_ay_b
}.
\]

\subsection*{Через длины и угол}

Если \(\varphi\) --- угол между ненулевыми векторами,
\[
0^\circ\le\varphi\le180^\circ,
\]
то
\[
\boxed{
\vect a\cdot\vect b
=
|\vect a|\,|\vect b|\cos\varphi
}.
\]

Следовательно,
\[
\boxed{
\displaystyle
\cos\varphi
=
\frac{x_ax_b+y_ay_b}
{\sqrt{x_a^2+y_a^2}\,\sqrt{x_b^2+y_b^2}}
},
\qquad
|\vect a|\ne0,
\quad
|\vect b|\ne0.
\]

Для ненулевых векторов полезен критерий перпендикулярности:
\[
\boxed{
\vect a\perp\vect b
\iff
\vect a\cdot\vect b=0
}.
\]

Также знак скалярного произведения позволяет быстро проверить характер угла:
\[
\vect a\cdot\vect b>0\Rightarrow\varphi<90^\circ,
\]
\[
\vect a\cdot\vect b=0\Rightarrow\varphi=90^\circ,
\]
\[
\vect a\cdot\vect b<0\Rightarrow\varphi>90^\circ.
\]

\textbf{Пример.}
Пусть
\[
\vect a=(1;2),
\qquad
\vect b=(-2;-1).
\]
Тогда
\[
\vect a\cdot\vect b
=1\cdot(-2)+2\cdot(-1)
=-4,
\]
\[
|\vect a|=\sqrt5,
\qquad
|\vect b|=\sqrt5.
\]
Поэтому
\[
\cos\varphi=-\frac45,
\]
а значит,
\[
\boxed{
\varphi=\arccos\!\left(-\frac45\right)
}.
\]
Отрицательный косинус согласуется с тем, что угол тупой.

\section*{6. Составные выражения: порядок действий}
\sectionrule

Если в скалярном произведении участвует сумма или разность векторов,
удобнее сначала выполнить операции с векторами.

Пусть
\[
\vect a=(-1;3),
\qquad
\vect b=(4;1),
\qquad
\vect c=(2;t),
\]
и известно, что
\[
(\vect a+\vect b)\cdot\vect c=0.
\]
Сначала
\[
\vect a+\vect b=(3;4).
\]
Затем
\[
(3;4)\cdot(2;t)=3\cdot2+4t.
\]
По условию
\[
6+4t=0,
\]
откуда
\[
\boxed{t=-\frac32}.
\]

\section*{7. Координатный метод в геометрии}
\sectionrule

Если координаты в условии отсутствуют, их часто выгодно \textbf{ввести самостоятельно}.
Используйте декартову систему координат: оси должны быть взаимно перпендикулярны,
а единичный масштаб по обеим осям --- одинаковым.

Удобное начало координат часто совпадает с:
\begin{itemize}[leftmargin=7mm,itemsep=1.5mm]
\checkitem вершиной фигуры;
\checkitem серединой стороны;
\checkitem центром симметрии или точкой пересечения диагоналей;
\checkitem основанием высоты.
\end{itemize}

\textbf{Пример: прямоугольник.}
Пусть
\[
AB=6,
\qquad
AD=8.
\]
Выберем
\[
A(0;0),\quad B(0;6),\quad D(8;0),\quad C(8;6).
\]
Тогда
\[
\overrightarrow{AB}=(0;6),
\qquad
\overrightarrow{AD}=(8;0),
\]
поэтому
\[
\overrightarrow{AB}+\overrightarrow{AD}=(8;6),
\]
и
\[
\left|
\overrightarrow{AB}+\overrightarrow{AD}
\right|
=\sqrt{8^2+6^2}=10.
\]
Тот же результат следует из правила параллелограмма:
\[
\overrightarrow{AB}+\overrightarrow{AD}=\overrightarrow{AC}.
\]

Пусть диагонали пересекаются в точке \(O\). Поскольку диагонали прямоугольника делятся точкой пересечения пополам,
\[
O(4;3).
\]
\[
\overrightarrow{AO}=(4;3),
\qquad
\overrightarrow{DO}=(-4;3).
\]
Следовательно,
\[
\overrightarrow{AO}+\overrightarrow{DO}=(0;6),
\]
и
\[
\boxed{
\left|
\overrightarrow{AO}+\overrightarrow{DO}
\right|=6
}.
\]

\begin{center}
\begin{tikzpicture}[scale=0.55]
\tkzDefPoint(0,0){A}
\tkzDefPoint(0,6){B}
\tkzDefPoint(8,6){C}
\tkzDefPoint(8,0){D}
\tkzInterLL(A,C)(B,D)
\tkzGetPoint{O}
\tkzDrawPolygon(A,B,C,D)
\tkzDrawSegments(A,C B,D)
\tkzDrawPoints(A,B,C,D,O)
\tkzLabelPoints[below left](A)
\tkzLabelPoints[above left](B)
\tkzLabelPoints[above right](C)
\tkzLabelPoints[below right](D)
\tkzLabelPoints[above](O)
\tkzLabelSegment[left](A,B){$6$}
\tkzLabelSegment[below](A,D){$8$}
\end{tikzpicture}
\end{center}

\clearpage
\section*{8. Полезная геометрическая заготовка: равносторонний треугольник}
\sectionrule

Для равностороннего треугольника со стороной \(a\) полезно помнить высоту
\[
\boxed{h=\frac{a\sqrt3}{2}}.
\]
Высота одновременно является медианой, поэтому по теореме Пифагора
\[
h^2+\left(\frac a2\right)^2=a^2.
\]
Отсюда
\[
h^2=\frac{3a^2}{4},
\qquad
h=\frac{a\sqrt3}{2}.
\]
В координатных задачах начало координат часто удобно помещать в основание такой высоты.

\section*{9. Сонаправленные и противоположно направленные векторы}
\sectionrule

Если ненулевые векторы \(\vect c\) и \(\vect b\) сонаправлены, то
\[
\boxed{
\vect c=k\vect b,
\qquad
k>0
}.
\]
Если они направлены противоположно, то
\[
\boxed{
\vect c=k\vect b,
\qquad
k<0
}.
\]
При умножении вектора на число его длина изменяется по правилу
\[
\boxed{|k\vect b|=|k|\,|\vect b|}.
\]

\textbf{Пример.}
Пусть
\[
\vect a=(-2;4),
\qquad
\vect b=(2;-1),
\]
вектор \(\vect c\) сонаправлен с \(\vect b\), причём
\[
|\vect c|=|\vect a|.
\]
Из сонаправленности
\[
\vect c=k\vect b=(2k;-k),
\qquad
k>0.
\]
Далее
\[
|\vect c|
=\sqrt{(2k)^2+(-k)^2}
=\sqrt{5k^2}
=|k|\sqrt5,
\]
а
\[
|\vect a|
=\sqrt{(-2)^2+4^2}
=\sqrt{20}
=2\sqrt5.
\]
Поэтому
\[
|k|\sqrt5=2\sqrt5,
\]
то есть
\[
|k|=2.
\]
Так как \(k>0\), получаем
\[
k=2,
\]
а значит,
\[
\boxed{\vect c=(4;-2)}.
\]

\section*{10. Типичные ошибки и самопроверка}
\sectionrule

\begin{itemize}[leftmargin=7mm,itemsep=2mm]
\checkitem В \(\overrightarrow{AB}\) перепутаны начало и конец. Проверка: всегда \(B-A\).
\checkitem При движении влево или вниз потерян знак «минус».
\checkitem В формуле длины одна из координат не возведена в квадрат.
\checkitem При поиске \(|\vect a|^2\) без необходимости извлекается корень.
\checkitem Смешиваются \(|\vect a|+|\vect b|\) и \(|\vect a+\vect b|\).
\checkitem В скалярном произведении перемножены несоответствующие координаты.
\checkitem При поиске угла найден только \(\cos\varphi\), хотя требуется сам \(\varphi\).
\checkitem Если \(\cos\varphi=c<0\), то \(\varphi=\arccos c\); дополнительная поправка не нужна.
\checkitem При сонаправленности допускается \(k<0\), хотя должно быть \(k>0\).
\checkitem В геометрии введены косые оси или разный единичный масштаб, после чего неверно применена формула длины.
\end{itemize}

\clearpage
\section*{11. Итоговый чек-лист формул}
\sectionrule

\begin{table}[h!]
\centering
\begin{tabularx}{\linewidth}{@{}>{\bfseries}p{0.36\linewidth}X@{}}
\toprule
Что требуется & Формула \\
\midrule
Координаты \(\overrightarrow{AB}\)
& \((x_B-x_A;\,y_B-y_A)\) \\[1.5mm]
Длина
& \(|\vect a|=\sqrt{x_a^2+y_a^2}\) \\[1.5mm]
Квадрат длины
& \(|\vect a|^2=x_a^2+y_a^2\) \\[1.5mm]
Умножение на число
& \(k\vect a=(kx_a;ky_a)\) \\[1.5mm]
Сумма
& \(\vect a+\vect b=(x_a+x_b;\,y_a+y_b)\) \\[1.5mm]
Разность
& \(\vect a-\vect b=(x_a-x_b;\,y_a-y_b)\) \\[1.5mm]
Скалярное произведение
& \(\vect a\cdot\vect b=x_ax_b+y_ay_b\) \\[1.5mm]
Скалярное произведение через угол
& \(\vect a\cdot\vect b=|\vect a||\vect b|\cos\varphi\) \\[1.5mm]
Косинус угла
& \(\displaystyle \cos\varphi=\frac{\vect a\cdot\vect b}{|\vect a||\vect b|}\) \\[2mm]
Перпендикулярность
& \(\vect a\perp\vect b\iff\vect a\cdot\vect b=0\) для ненулевых векторов \\[1.5mm]
Сонаправленность
& \(\vect c=k\vect b,\ k>0\) \\
\bottomrule
\end{tabularx}
\end{table}

% =========================================================
% Блок-схема
% =========================================================

\clearpage
\thispagestyle{fancy}

\begin{center}
{\Large\bfseries\textcolor{darkaccent}{Алгоритм: координатный метод в задачах на векторы}}
\end{center}

\vspace{5mm}

\begin{center}
\begin{tikzpicture}[node distance=7mm and 12mm]

\node[flowstart] (start) {Начало решения};

\node[flowdecision,below=of start] (given)
{Координаты\\уже заданы?};

\node[flowaction,below left=8mm and 6mm of given,text width=47mm] (choose)
{Введите декартову систему координат. Выберите удобное начало: вершина, середина, центр симметрии или основание высоты.};

\node[flowaction,below right=8mm and 6mm of given,text width=47mm] (read)
{Аккуратно выпишите данные координаты точек и векторов.};

\node[flowaction,below=39mm of given,text width=67mm] (vecs)
{Найдите координаты нужных векторов: конец минус начало.};

\node[flowaction,below=of vecs,text width=67mm] (ops)
{Выполните требуемое векторное выражение. При необходимости найдите его длину.};

\node[flowdecision,below=of ops] (angle)
{Нужно найти угол?};

\node[flowaction,below right=8mm and 6mm of angle,text width=47mm] (cosine)
{Вычислите \(\vect a\cdot\vect b\), длины векторов, затем \(\cos\varphi\) и \(\varphi\).};

\node[flowaction,below=39mm of angle,text width=67mm] (check)
{Проверьте знаки, порядок координат и геометрический смысл ответа.};

\node[flowstart,below=of check] (finish) {Ответ};

\draw[flowarrow] (start) -- (given);

\draw[flowarrow] (given) -- node[above left,font=\small]{нет} (choose);
\draw[flowarrow] (given) -- node[above right,font=\small]{да} (read);

\draw[flowarrow] (choose.south) |- (vecs.west);
\draw[flowarrow] (read.south) |- (vecs.east);

\draw[flowarrow] (vecs) -- (ops);
\draw[flowarrow] (ops) -- (angle);

\draw[flowarrow] (angle) -- node[above right,font=\small]{да} (cosine);
\draw[flowarrow] (angle.west) -- ++(-18mm,0) |- node[pos=0.18,left,font=\small]{нет} (check.west);
\draw[flowarrow] (cosine.south) |- (check.east);

\draw[flowarrow] (check) -- (finish);

\end{tikzpicture}
\end{center}

\clearpage

% =========================================================
% Тренировочный блок
% =========================================================

\section*{Тренировочный блок --- Путь А}
\sectionrule

\textbf{Инструкция.}
Решите задания самостоятельно. Промежуточные координаты рекомендуется выписывать полностью:
это снижает риск ошибок в знаках и порядке вычитания.

\begin{enumerate}[leftmargin=8mm,label=\textbf{\arabic*.},itemsep=8mm]

\item
Даны точки
\[
A(-2;5),
\qquad
B(4;1).
\]
Найдите координаты вектора \(\overrightarrow{AB}\) и его длину.

\item
Даны векторы
\[
\vect a=(3;-1),
\qquad
\vect b=(1;2).
\]
Найдите
\[
2\vect a-\vect b
\]
и длину полученного вектора.

\item
Даны векторы
\[
\vect a=(1;1),
\qquad
\vect b=(2;0).
\]
Найдите угол между этими векторами.

\item
В прямоугольнике \(ABCD\)
\[
AB=8,
\qquad
AD=12.
\]
Диагонали пересекаются в точке \(O\).
Найдите
\[
\left|
\overrightarrow{AO}+\overrightarrow{DO}
\right|.
\]
Рекомендуется использовать координатный метод.

\item
Даны
\[
\vect a=(-6;8),
\qquad
\vect b=(3;-4).
\]
Вектор \(\vect c\) сонаправлен с \(\vect b\), причём
\[
|\vect c|=|\vect a|.
\]
Найдите координаты \(\vect c\) и значение
\[
x_c+y_c.
\]

\end{enumerate}

% =========================================================
% Ответы
% =========================================================

\clearpage

\section*{Ответы}
\sectionrule

\begin{table}[h!]
\centering
\begin{tabularx}{\linewidth}{@{}p{0.12\linewidth}X@{}}
\toprule
\textbf{№} & \textbf{Ответ} \\
\midrule
1
&
\(\overrightarrow{AB}=(6;-4),\quad |\overrightarrow{AB}|=2\sqrt{13}.\)
\\[2mm]
2
&
\(2\vect a-\vect b=(5;-4),\quad |2\vect a-\vect b|=\sqrt{41}.\)
\\[2mm]
3
&
\(45^\circ.\)
\\[2mm]
4
&
\(8.\)
\\[2mm]
5
&
\(\vect c=(6;-8),\quad x_c+y_c=-2.\)
\\
\bottomrule
\end{tabularx}
\end{table}

\vfill

\begin{center}
\small\textcolor{textgray}{Лёвин Артём Александрович эксклюзивно для Матвея Горбачева}
\end{center}

\end{document}